% Une ligne commentaire débute par le caractère « % »

\documentclass[a4paper]{article}

% Options possibles : 10pt, 11pt, 12pt (taille de la fonte)
%                     oneside, twoside (recto simple, recto-verso)
%                     draft, final (stade de développement)

\usepackage[utf8]{inputenc}   % LaTeX, comprends les accents !
\usepackage[T1]{fontenc}      % Police contenant les caractères français
\usepackage[francais]{babel}  
\usepackage{fullpage}
\usepackage[colorlinks=true, linkcolor=blue]{hyperref}


\usepackage{graphicx}  % pour inclure des images

%\pagestyle{headings}        % Pour mettre des entêtes avec les titres
                              % des sections en haut de page

 \title{  Simulation de fluides en 2D\\         % Les paramètres du titre : titre, auteur, date
  Projet de programmation 2}          
\author{ PEYREROL Axel, MILEZI Gaëlle, BARRAUX Satine, DELRIEU Saurane\\
  \emph{L3 informatique}\\
  Faculté des Sciences\\
Université de Montpellier.}
\date{2025-2026}             



\begin{document}

\maketitle                    % Faire un titre utilisant les données
                              % passées à \title, \author et \date

%\begin{center}               % pour centrer 
 % \includegraphics[scale=1]{}   % insertion d'une image
%\end{center}

\begin{abstract}     % Résumé du travail

  \emph{Description très succinte du problème et des différentes étapes de réalisation}

\end{abstract}
\newpage
%\dominitoc  % initializer les minitoc
\tableofcontents
\newpage
\section*{Introduction} % Commencer une section, * permet de ne pas numéroté l'introduction
\emph{1-2 pages}
	%Présentation du Sujet (1-2 pages) 
\begin{itemize}%pour faire une liste à puces
\item Présenter le problème étudié et le contexte dans lequel le projet se positionne.\\
   Le but de ce projet est de simuler un fluide en 2D. 
  Pour cela, nous avons choist d'implementer la methode
  de Jos Stam, "Stable Fluids", qui permet de simuler un 
  fluide de manière stable et rapide.
\item Motiver l'intérêt du problème étudié par rapport à votre parcours d'études et au monde de l'informatique.\\
    La simulation des fluides est un sujet important dans diverse domaine que cesoit
    en biologie a travers la simulation de flux sanguin, en physique pour la simulation de l'air 
    ou dans l'ingénierie.il est est utilisé dans de nombreux domaines, faisant ainsi de ce sujet
    un sujet d'étude intéressant pour les étudiants en informatique.
    choisir ce sujet etait pour nous une evidence car il nous permet de mette 
    en avant nos connaissance en informatique mais egalement celle aquit durant 
    notre licence de mathematiques tout en y apportant de nouvelle conaisance
    et une plus grande ouverture d'esprit grace a ces nobreuse notion de phisique.
\item Donner le cahier des charges détaillé
\end{itemize}


\section{Développements Logiciel : Conception, Modélisation, Implémentation} 
\emph{8 à 15 pages}

Il se peut que la conception du logiciel (3) et l'analyse algorithmique (4) aient un poids différent dans votre projet. Les points décrits ci-dessous seront ainsi plus ou moins développés en fonction du travail réalisé en accord avec l'encadrant. 
\begin{enumerate}%pour faire une liste numérotée
\item Présenter les développements logiciel réalisés dans le cadre du projet.\\
\item Présenter les principaux modules du logiciel développé dans le cadre du projet. Utiliser le langage UML pour la modélisation : donner le diagramme de cas d'utilisation et le diagramme des classes.
\item Décrire les fonctionnalités de l'interface graphique implémentée (si votre logiciel dispose d'une interface graphique).
\item Décrire le format des données en entrée ou encore les conventions utilisées pour les entrées de vos programmes. Décrire les procédures de lecture et validation des entrées.
\item Statistiques : nombre de modules/composantes/classes/scripts développés. Nombre de lignes de code.
\end{enumerate}

\section{Algorithmes et Structures de Données}
\emph{2 à 3 pages}
\begin{enumerate}%pour faire une liste numérotée
\item 	Présenter les principales structures de données définies dans le cadre du projet.\\
      Dans le cadre de notre projet, nous avons défini plusieurs structures de données pour représenter les différentes composantes
      Ces diefferntes composantes sont ranger dans trois principales classes : 
      la classe "main", la classe "fluidRender" et la classe "display".
      Dans la classe "main" on retrouve la bases du code ainsi que l'affichage
      de celui ci.
      La classe "fluidRender" contient notament les fonctions de calcul du fluides
      ainsi que celle delimitant les bort de ceui ci.
      Enfin la classe "display" contient les fonctions d'affichage du fluides ou encore 
      des differents obstacles.
\item 	Présenter les principaux algorithmes implémentés. En illustrer le fonctionnement avec des exemples. Attention : il s'agit ici de choisir 1 ou 2 algorithmes intéressants, et non pas de présenter tous les algorithmes implémentés.
\item 	Évaluer la complexité théorique en temps des algorithmes présentés
\end{enumerate}

\section{Analyse des résultats }
\emph{2 à 4 pages}
L'analyse expérimentale sera faites dans cette section.   Elle sera plus ou moins importante dans votre projet comme pour les deux points précédent, cette section devra donc être développée en conséquence.

\begin{enumerate}%pour faire une liste numérotée
\item 	Illustrer les performances ainsi que l'efficacité du logiciel implémenté à l'aide de graphiques.
\item Analyser (et comparer, si plusieurs) les performances des solutions implémentées.
\item	Présenter les bancs d'essais (ou les procédures utilisées pour la génération des données) utilisés pour les tests du logiciel.
\end{enumerate}

\section{Gestion du Projet}
\emph{ (1-2 pages)}
\begin{enumerate}%pour faire une liste numérotée
\item 	Présenter la gestion du projet et les documents de planification rédigés (par exemple, le diagramme de Gantt).
\item	Discuter les changements majeurs effectués en cours de projet.\\
      Durant le projet, nous avons du faire face a plusieurs changement majeur notamment 
      dans la partie de l'implémentation du projet. En effet, 
      nous avions commencer ce project en implementant une grille ainsi que
      diverse fonction telle que celle afiichant les vecteurs
      mais au moment de l'implementation du fluides naus avaons du tout revoir
      car les premier essaie ne nous permetait pas d'obtenir le resultat escompté. 
      Nous avons ainsi du revoir notre code en implementant la methode de Jos Stam.
      \\
      \\
      Une fois cette methode implementé nous avons decider de nous partager diverse 
      tache afin d'etre des plus efficace. Ainsi, chacun d'etre nous a rencotré des
      difficulter.
      pour les creation d'obstacle, il a fallut dans un premier temps readapter la fonction
      qui delimiter les bord du fluides afin que celui ci prenne en compte l'object
      ce qui a crer de nombreux probleme (la disparition du fluide a travers l'obstacle ou la fuite du fluide a travers les bords).
      ou encore des probleme de dependance entre les diferents obstacles lorsque que nous avons voulu
      ajouter plusieurs obstacles mobile. De plus, l'implementation d'obstacle ayant une autre forme
      que la boule pose de nombreux probleme notamment dans la partie de delimitation du fluides ou encore dans la partie d'affichage.
      Pour resoudre ces probleme il a fallut revoir le code en changement la facon de voir 
      les obstacles mais également en les implementant de maniere differente.
      Ainsi, les obstacle ont ete ranger une une liste afin que le premier obstacle soit reconnaisable
      et ceux ci sont ainsi selectionner a partir du premier de la liste lorsqu'il fuat
      les afficher et c'est le dernier obstacles qui est effacer lorsque l'on desire e retirer un.

\end{enumerate}

\section{Bilan et Conclusions}

	Indiquer les fonctionnalités mises en œuvre par rapport au cahier des charges de départ, les points ouvertes.
	Les perspectives du projet décrivent ce que l’on pourrait ajouter, modifier, réécrire afin d’améliorer le projet. 
	
\begin{thebibliography}{9}
\bibitem{texbook}
Donald E. Knuth (1986) \emph{The \TeX{} Book}, Addison-Wesley Professional.

\bibitem{lamport94}
Leslie Lamport (1994) \emph{\LaTeX: a document preparation system}, Addison
Wesley, Massachusetts, 2nd ed.
\end{thebibliography}

\end{document}

